\chapter{Échantillonnage et simulation}
\section*{Exercice 23.1 -- Moyenne d'échantillon}
Une variable a moyenne $8$ et écart type $6$. Pour $n=100$ observations indépendantes de même loi, donner les moments de leur moyenne.

\section*{Exercice 23.2 -- Taille et précision}
Combien faut-il multiplier la taille d'un échantillon indépendant pour diviser par $3$ l'écart type de sa moyenne ?

\section*{Exercice 23.3 -- Tchebychev}
Pour une moyenne de $n$ Bernoulli indépendantes de paramètre $1/2$, majorer la probabilité d'un écart d'au moins $0{,}1$ à $1/2$. Quel $n$ rend la borne au plus $0{,}05$ ?

\section*{Exercice 23.4 -- Deux écarts types}
Pour une Bernoulli de paramètre $0{,}1$ et $n=1$, calculer la probabilité d'être dans $\mu\pm2\sigma$.

\section*{Exercice 23.5 -- Monte-Carlo}
Sur $10\,000$ points simulés dans le carré unité, $7\,860$ sont dans le quart de disque. Donner l'estimation de $\pi$ et une borne théorique de son écart type.

\section*{Exercice 23.6 -- Pourquoi l'indépendance compte}
Si toutes les observations sont identiques à la même variable $X$, quelle est la variance de leur moyenne ?
